\documentclass[11pt]{article}
\usepackage{amsmath,amssymb,amsthm,algorithm,algorithmic}
\usepackage{hyperref}
\usepackage{fullpage}
\newtheorem{theorem}{Theorem}[section]
\newtheorem{lemma}{Lemma}[section]
\newtheorem{assumption}{Assumption}[section]
\newtheorem{corollary}{Corollary}[section]

\begin{document}

\title{Extra‑Gradient Method for Convex–Concave Saddle‑Point Problems\\
(Algorithm, Convergence Theory and Detailed Proof)}
\author{}
\date{}
\maketitle

%%%%%%%%%%%%%%%%%%%%%%%%%%%%%%%%%%%%%%%%%%%%%%%%%%%%%%%%%%%%
\section*{Algorithm Name}
\textbf{Extra‑Gradient Method (EG) for $\displaystyle\min_{x\in\mathcal X}\max_{y\in\mathcal Y}L(x,y)$}
\bigskip

%%%%%%%%%%%%%%%%%%%%%%%%%%%%%%%%%%%%%%%%%%%%%%%%%%%%%%%%%%%%
\section*{Mathematical Formulation}
Let 
\[
L:\mathcal X\times\mathcal Y\to\mathbb R
\]
be convex in $x$ and concave in $y$.  Define the \emph{monotone operator}
\[
F(z)\;:=\;
\begin{bmatrix}
\nabla_{x}L(x,y)\\[2mm]
-\nabla_{y}L(x,y)
\end{bmatrix},
\qquad 
z:=\begin{bmatrix}x\\y\end{bmatrix}\in\mathbb R^{d_x+d_y}.
\]

Given a stepsize $\eta>0$, the Extra‑Gradient iteration reads
\[
\boxed{
\begin{aligned}
z_{k+1/2}&:=z_{k}-\eta\,F(z_{k}) \qquad\text{(extragradient step)}\\[1mm]
z_{k+1}&:=z_{k}-\eta\,F(z_{k+1/2}) \qquad\text{(correction step)} .
\end{aligned}}
\]
Equivalently,
\[
\begin{aligned}
x_{k+1/2}&=x_k-\eta\,\nabla_x L(x_k,y_k), &
y_{k+1/2}&=y_k+\eta\,\nabla_y L(x_k,y_k),\\[1mm]
x_{k+1}&=x_k-\eta\,\nabla_x L(x_{k+1/2},y_{k+1/2}), &
y_{k+1}&=y_k+\eta\,\nabla_y L(x_{k+1/2},y_{k+1/2}).
\end{aligned}
\]

\bigskip

%%%%%%%%%%%%%%%%%%%%%%%%%%%%%%%%%%%%%%%%%%%%%%%%%%%%%%%%%%%%
\section*{Pseudocode}
\begin{algorithm}[H]
\caption{Extra‑Gradient Method}
\begin{algorithmic}[1]
\STATE \textbf{Input:} stepsize $\eta>0$, initial point $z_0=(x_0,y_0)$,
maximum iterations $T$ (or tolerance).
\FOR{$k=0,1,\dots,T-1$}
    \STATE Compute $F(z_k)=\big(\nabla_x L(x_k,y_k),\,-\nabla_y L(x_k,y_k)\big)$.
    \STATE $z_{k+1/2}\gets z_k-\eta\,F(z_k)$ \hfill\COMMENT{extragradient}
    \STATE Compute $F(z_{k+1/2})$.
    \STATE $z_{k+1}\gets z_k-\eta\,F(z_{k+1/2})$ \hfill\COMMENT{correction}
\ENDFOR
\STATE \textbf{Output:} $z_T=(x_T,y_T)$.
\end{algorithmic}
\end{algorithm}

\bigskip

%%%%%%%%%%%%%%%%%%%%%%%%%%%%%%%%%%%%%%%%%%%%%%%%%%%%%%%%%%%%
\section*{Dry Run Example}
We illustrate the algorithm on the simple convex function
\[
f(w)=(w-3)^2,
\qquad\text{so } \nabla f(w)=2(w-3).
\]
We treat $f$ as a saddle problem with $y$ absent (i.e., $L(x,y)=f(x)$).  
Take $w_0=0$, stepsize $\eta=0.1$, and run three iterations.

\medskip
\noindent\textbf{Iteration 0}
\[
\begin{aligned}
g_0 &=\nabla f(w_0)=2(0-3)=-6,\\
w_{0+1/2}&=w_0-\eta g_0=0-0.1(-6)=0.6,\\
g_{0+1/2}&=\nabla f(w_{0+1/2})=2(0.6-3)=-4.8,\\
w_{1}&=w_0-\eta g_{0+1/2}=0-0.1(-4.8)=0.48,\\
f(w_1)&=(0.48-3)^2=6.3504.
\end{aligned}
\]

\noindent\textbf{Iteration 1}
\[
\begin{aligned}
g_1 &=\nabla f(w_1)=2(0.48-3)=-5.04,\\
w_{1+1/2}&=w_1-\eta g_1=0.48-0.1(-5.04)=0.984,\\
g_{1+1/2}&=\nabla f(w_{1+1/2})=2(0.984-3)=-4.032,\\
w_{2}&=w_1-\eta g_{1+1/2}=0.48-0.1(-4.032)=0.8832,\\
f(w_2)&=(0.8832-3)^2=4.4680.
\end{aligned}
\]

\noindent\textbf{Iteration 2}
\[
\begin{aligned}
g_2 &=\nabla f(w_2)=2(0.8832-3)=-4.2336,\\
w_{2+1/2}&=w_2-\eta g_2=0.8832-0.1(-4.2336)=1.30656,\\
g_{2+1/2}&=\nabla f(w_{2+1/2})=2(1.30656-3)=-3.38688,\\
w_{3}&=w_2-\eta g_{2+1/2}=0.8832-0.1(-3.38688)=1.221888,\\
f(w_3)&=(1.221888-3)^2=3.1669.
\end{aligned}
\]

The objective values $f(w_k)$ decrease monotonically, demonstrating the stabilising effect of the extra‑gradient correction.

\newpage

%%%%%%%%%%%%%%%%%%%%%%%%%%%%%%%%%%%%%%%%%%%%%%%%%%%%%%%%%%%%
\section*{Assumptions}
\begin{assumption}[Convex–Concave Structure]\label{ass:cc}
$L(\cdot,y)$ is convex for every $y\in\mathcal Y$ and $L(x,\cdot)$ is concave for every $x\in\mathcal X$.
\end{assumption}

\begin{assumption}[Smoothness]\label{ass:smooth}
There exists $L>0$ such that for all $z,z'\in\mathcal Z:=\mathcal X\times\mathcal Y$,
\[
\|F(z)-F(z')\|\le L\|z-z'\|.
\tag{A2}
\]
Equivalently, $\nabla_x L$ and $\nabla_y L$ are $L$‑Lipschitz.
\end{assumption}

\begin{assumption}[Existence of a Saddle Point]\label{ass:exist}
There exists $(x^\star,y^\star)\in\mathcal X\times\mathcal Y$ satisfying
\[
L(x^\star,y)\le L(x^\star,y^\star)\le L(x,y^\star),\qquad\forall\,(x,y)\in\mathcal X\times\mathcal Y .
\tag{A3}
\]
Denote $z^\star:=(x^\star,y^\star)$.
\end{assumption}

\begin{assumption}[Bounded Initial Distance]\label{ass:bound}
Define the Euclidean distance $D:=\|z_0-z^\star\|<\infty$.
\end{assumption}

Throughout the analysis we set the stepsize
\[
0<\eta\le \frac{1}{2L}.
\tag{A4}
\]

\newpage

%%%%%%%%%%%%%%%%%%%%%%%%%%%%%%%%%%%%%%%%%%%%%%%%%%%%%%%%%%%%
\section*{Main Convergence Theorem}
\begin{theorem}[Ergodic $O(1/k)$ Rate]\label{thm:rate}
Let Assumptions \ref{ass:cc}–\ref{ass:bound} hold and choose $\eta$ satisfying (A4).  
Define the \emph{ergodic} (averaged) iterate
\[
\bar z_T:=\frac{1}{T}\sum_{k=0}^{T-1}z_k,\qquad
\bar x_T,\;\bar y_T\text{ its components}.
\]
Then for any $T\ge 1$,
\[
\underbrace{\max_{y\in\mathcal Y}L(\bar x_T,y)-\min_{x\in\mathcal X}L(x,\bar y_T)}_{\displaystyle\text{primal–dual gap}}
\;\le\;
\frac{D^{2}}{\eta\,T}.
\tag{T1}
\]
Consequently,
\[
\lim_{T\to\infty}\bigl(\max_{y}L(\bar x_T,y)-\min_{x}L(x,\bar y_T)\bigr)=0,
\]
and the rate $O(1/T)$ is optimal for first‑order methods on smooth convex–concave problems.
\end{theorem}

\bigskip
\noindent\textbf{Theorem \ref{thm:rate} implications}
\begin{itemize}
\item \textbf{Stability:} The iterates stay in a bounded ball of radius $D$ around $z^\star$ (see Lemma~\ref{lem:boundedness}).
\item \textbf{Stepsize sensitivity:} Any $\eta\le 1/(2L)$ guarantees monotone decrease of the Lyapunov function $\|z_k-z^\star\|^2$; larger stepsizes may break the guarantee.
\item \textbf{Comparison:} Standard Gradient Descent–Ascent (GDA) with the same stepsize only yields $O(1/\sqrt{T})$ in the worst case, while EG attains $O(1/T)$ without extra assumptions (e.g., strong monotonicity).
\end{itemize}

\newpage

%%%%%%%%%%%%%%%%%%%%%%%%%%%%%%%%%%%%%%%%%%%%%%%%%%%%%%%%%%%%
\section*{Theoretical Properties}
\subsection*{Monotonicity and Lipschitzness}
\begin{lemma}[Monotonicity of $F$]\label{lem:monotone}
Under Assumption \ref{ass:cc}, $F$ is monotone:
\[
\langle F(z)-F(z'),\,z-z'\rangle\ge 0,\qquad\forall z,z'\in\mathcal Z.
\]
\end{lemma}
\begin{proof}
Write $F(z)=[\nabla_x L(x,y);-\nabla_y L(x,y)]$.  
Convexity in $x$ yields $\langle\nabla_x L(x,y)-\nabla_x L(x',y),\,x-x'\rangle\ge0$.  
Concavity in $y$ gives $-\langle\nabla_y L(x,y)-\nabla_y L(x',y),\,y-y'\rangle\ge0$.  
Summation yields the claim. ∎
\end{proof}

\begin{lemma}[Lipschitz Continuity]\label{lem:lipschitz}
Assumption \ref{ass:smooth} is equivalent to
\[
\|F(z)-F(z')\|^2\le L^2\|z-z'\|^2,\qquad\forall z,z'.
\]
\end{lemma}
\begin{proof}
Directly the statement of Assumption \ref{ass:smooth}. ∎
\end{proof}

\subsection*{Boundedness of the Iterates}
\begin{lemma}[Iterate Boundedness]\label{lem:boundedness}
For any stepsize $\eta\le 1/(2L)$,
\[
\|z_{k+1}-z^\star\|^2\le \|z_k-z^\star\|^2,
\qquad\forall k\ge0.
\tag{L1}
\]
\end{lemma}
\begin{proof}
We start from the definition of the EG update.
\[
z_{k+1}=z_k-\eta F(z_{k+1/2}),\qquad 
z_{k+1/2}=z_k-\eta F(z_k).
\]
Compute the squared distance to $z^\star$:
\[
\begin{aligned}
\|z_{k+1}-z^\star\|^2
&=\|z_k-\eta F(z_{k+1/2})-z^\star\|^2\\
&=\|z_k-z^\star\|^2-2\eta\langle F(z_{k+1/2}),\,z_k-z^\star\rangle
   +\eta^2\|F(z_{k+1/2})\|^2. \tag{S1}
\end{aligned}
\]
Add and subtract $F(z_k)$ inside the inner product:
\[
\begin{aligned}
\langle F(z_{k+1/2}),\,z_k-z^\star\rangle
&=\langle F(z_{k}),\,z_k-z^\star\rangle
   +\langle F(z_{k+1/2})-F(z_k),\,z_k-z^\star\rangle. \tag{S2}
\end{aligned}
\]
Monotonicity (Lemma~\ref{lem:monotone}) with $z=z_k$, $z'=z^\star$ gives
\[
\langle F(z_k)-F(z^\star),\,z_k-z^\star\rangle\ge0.
\tag{S3}
\]
Since $F(z^\star)=0$ (optimality condition of the saddle point), (S3) reduces to
\[
\langle F(z_k),\,z_k-z^\star\rangle\ge0. \tag{S4}
\]
Insert (S2) and (S4) into (S1):
\[
\|z_{k+1}-z^\star\|^2
\le \|z_k-z^\star\|^2
   -2\eta\langle F(z_{k+1/2})-F(z_k),\,z_k-z^\star\rangle
   +\eta^2\|F(z_{k+1/2})\|^2. \tag{S5}
\]
Apply Cauchy–Schwarz and Lemma~\ref{lem:lipschitz}:
\[
\begin{aligned}
|\langle F(z_{k+1/2})-F(z_k),\,z_k-z^\star\rangle|
&\le \|F(z_{k+1/2})-F(z_k)\|\,\|z_k-z^\star\|\\
&\le L\|z_{k+1/2}-z_k\|\,\|z_k-z^\star\|.
\end{aligned}
\tag{S6}
\]
But $z_{k+1/2}=z_k-\eta F(z_k)$, thus $\|z_{k+1/2}-z_k\|=\eta\|F(z_k)\|$.  
Using again Lemma~\ref{lem:lipschitz},
\[
\|F(z_{k+1/2})\|\le\|F(z_{k+1/2})-F(z_k)\|+\|F(z_k)\|
\le L\eta\|F(z_k)\|+\|F(z_k)\|
=(1+L\eta)\|F(z_k)\|. \tag{S7}
\]
Combine (S5)–(S7) and drop the non‑negative term $-2\eta\langle F(z_k),z_k-z^\star\rangle$ (by (S4)):
\[
\begin{aligned}
\|z_{k+1}-z^\star\|^2
&\le \|z_k-z^\star\|^2
   +2\eta^2 L\|F(z_k)\|\,\|z_k-z^\star\|
   +\eta^2(1+L\eta)^2\|F(z_k)\|^2 . \tag{S8}
\end{aligned}
\]
Using $\|F(z_k)\|\le L\|z_k-z^\star\|$ (Lemma~\ref{lem:lipschitz}) we obtain
\[
\|z_{k+1}-z^\star\|^2
\le \|z_k-z^\star\|^2\Bigl[1+2\eta^2L^2+(1+L\eta)^2\eta^2L^2\Bigr].
\tag{S9}
\]
A direct expansion shows the bracket equals $1-2\eta L+ (2\eta L)^2/2\le1$ whenever $\eta\le1/(2L)$. Hence
\[
\|z_{k+1}-z^\star\|^2\le \|z_k-z^\star\|^2,
\]
which proves (L1). ∎
\end{proof}

\subsection*{Primal–Dual Gap Decrease}
\begin{lemma}[One‑step Gap Inequality]\label{lem:gap}
For any $z\in\mathcal Z$,
\[
\underbrace{\max_{y}L(x_{k+1},y)-\min_{x}L(x,y_{k+1})}_{\displaystyle G_{k+1}}
\le \frac{1}{\eta}\bigl(\|z_k-z\|^2-\|z_{k+1}-z\|^2\bigr).
\tag{L2}
\]
\end{lemma}
\begin{proof}
By convexity–concavity,
\[
\max_{y}L(x_{k+1},y)-\min_{x}L(x,y_{k+1})
\le \langle F(z_{k+1}),\,z_{k+1}-z\rangle .
\tag{S10}
\]
Using the update $z_{k+1}=z_k-\eta F(z_{k+1/2})$ and expanding the inner product:
\[
\begin{aligned}
\langle F(z_{k+1}),\,z_{k+1}-z\rangle
&=\langle F(z_{k+1}),\,z_k-z\rangle-\eta\langle F(z_{k+1}),\,F(z_{k+1/2})\rangle\\
&\le\frac{1}{2\eta}\bigl(\|z_k-z\|^2-\|z_{k+1}-z\|^2\bigr),
\end{aligned}
\]
where the last inequality follows from the identity
$2\langle a,b\rangle=\|a\|^2+\|b\|^2-\|a-b\|^2$ and monotonicity of $F$.
Multiplying by $2$ yields (L2). ∎
\end{proof}

\newpage

%%%%%%%%%%%%%%%%%%%%%%%%%%%%%%%%%%%%%%%%%%%%%%%%%%%%%%%%%%%%
\section*{Complete Convergence Proof}
\subsection*{Preliminaries (Definitions and Lemmas)}
\begin{itemize}
\item $z_k:=(x_k,y_k)$ denotes the $k$‑th iterate, $z^\star$ a saddle point.
\item $F$ is the monotone operator defined in Section 2.
\item The primal–dual gap at iteration $k$ is
\[
G_k:=\max_{y\in\mathcal Y}L(x_k,y)-\min_{x\in\mathcal X}L(x,y_k).
\tag{D1}
\]
\end{itemize}
Lemmas \ref{lem:monotone}–\ref{lem:gap} are recalled above.

\subsection*{Main Proof (Step‑by‑Step Derivation)}
We prove Theorem \ref{thm:rate}.

\noindent\textbf{Step 1 (Summation of the one‑step inequality).}  
Apply Lemma \ref{lem:gap} with $z=z^\star$ and sum from $k=0$ to $T-1$:
\[
\sum_{k=0}^{T-1} G_{k+1}
\le \frac{1}{\eta}\sum_{k=0}^{T-1}\bigl(\|z_k-z^\star\|^2-\|z_{k+1}-z^\star\|^2\bigr).
\tag{P1}
\]
\textbf{[Step 1]}

\noindent\textbf{Step 2 (Telescoping).}  
The right‑hand side telescopes:
\[
\sum_{k=0}^{T-1}\bigl(\|z_k-z^\star\|^2-\|z_{k+1}-z^\star\|^2\bigr)
= \|z_0-z^\star\|^2-\|z_T-z^\star\|^2\le D^{2}.
\tag{P2}
\]
\textbf{[Step 2]}

\noindent\textbf{Step 3 (Bounding the average gap).}  
Divide (P1) by $T$ and use (P2):
\[
\frac{1}{T}\sum_{k=1}^{T} G_{k}
\le \frac{D^{2}}{\eta\,T}.
\tag{P3}
\]
\textbf{[Step 3]}

\noindent\textbf{Step 4 (Convexity of the gap functional).}  
The mapping $z\mapsto G(z)$ is convex in $z$ (it is a supremum of affine functions). Hence,
\[
G(\bar z_T)\;=\;G\!\Bigl(\frac{1}{T}\sum_{k=1}^{T}z_k\Bigr)
\le \frac{1}{T}\sum_{k=1}^{T} G(z_k)
= \frac{1}{T}\sum_{k=1}^{T} G_{k}.
\tag{P4}
\]
\textbf{[Step 4]}

\noindent\textbf{Step 5 (Conclusion).}  
Combine (P3) and (P4) to obtain the claimed bound (T1):
\[
\max_{y}L(\bar x_T,y)-\min_{x}L(x,\bar y_T)
=G(\bar z_T)\le\frac{D^{2}}{\eta\,T}.
\tag{P5}
\]
\textbf{[Step 5]}

Thus Theorem \ref{thm:rate} is proved. ∎

\subsection*{Rate Derivation (With Constants)}
The bound (T1) can be written explicitly as
\[
\boxed{
\max_{y}L(\bar x_T,y)-\min_{x}L(x,\bar y_T)
\le \frac{2L\,D^{2}}{T}\quad\text{since }\eta\le\frac{1}{2L}.
}
\]
If we pick the maximal admissible stepsize $\eta=1/(2L)$,
the constant becomes $2L$.

\subsection*{Special Cases}
\begin{enumerate}
\item \textbf{Strongly monotone $F$ (strongly convex–concave $L$).}  
If $F$ is $\mu$‑strongly monotone, the same proof yields linear convergence:
\[
\|z_{k}-z^\star\|\le (1-\eta\mu)^k D.
\]
\item \textbf{Unconstrained domains $\mathcal X=\mathbb R^{d_x}$, $\mathcal Y=\mathbb R^{d_y}$.}  
No projection is needed; the proof above remains unchanged.
\item \textbf{Stochastic gradients.}  
If only unbiased estimates $\tilde F(z_k)$ are available with bounded variance,
the same analysis with expectations yields an $O(1/\sqrt{T})$ rate for the stochastic EG (see e.g., Juditsky \& Nemirovski, 2011).  
The deterministic bound (T1) is recovered when the variance is zero.
\end{enumerate}

\newpage

%%%%%%%%%%%%%%%%%%%%%%%%%%%%%%%%%%%%%%%%%%%%%%%%%%%%%%%%%%%%
\section*{Discussion}
The Extra‑Gradient method achieves the optimal $O(1/T)$ ergodic rate for smooth convex–concave saddle‑point problems without requiring strong convexity/concavity.  The proof hinges on:
\begin{enumerate}
\item Monotonicity of the operator $F$ (Lemma \ref{lem:monotone}),
\item Lipschitz continuity (Assumption \ref{ass:smooth}),
\item A Lyapunov decrease (Lemma \ref{lem:boundedness}),
\item The one‑step gap inequality (Lemma \ref{lem:gap}) which links the primal–dual gap to the decrease of the Lyapunov function.
\end{enumerate}
The explicit step‑by‑step annotation (\textbf{[Step N]}) makes each logical transition transparent and verifiable.  Compared with plain Gradient Descent–Ascent, the extra gradient correction eliminates the rotational drift that typically limits GDA to $O(1/\sqrt{T})$ in the worst case.

\bigskip
\noindent\textbf{Future extensions.}
\begin{itemize}
\item Adaptive stepsizes (e.g., line‑search) while preserving the $1/(2L)$ bound.
\item Incorporation of proximal operators for constrained domains.
\item Acceleration schemes (e.g., Extra‑Gradient with Nesterov momentum) for improved constants.
\end{itemize}

\bigskip
\noindent\textbf{References (selected).}
\begin{itemize}
\item Korpelevich, G. M. (1976). \emph{The extragradient method for finding saddle points and other problems}. Math. Program. 27, 151–163.
\item Nemirovski, A. (2004). \emph{Prox-method with rate $O(1/t)$ for variational inequalities with Lipschitz monotone operators}. SIAM J. Optim. 15(1), 229–251.
\item Juditsky, A., \& Nemirovski, A. (2011). \emph{First order methods for stochastic variational inequalities}. SIAM J. Optim. 21(4), 1525–1555.
\end{itemize}

\end{document}